\subsection{Interplay between Entanglement and Subsystem Fermion Parity}\label{sec:EE-calculation-odd}

A subtlety arises when the number of lattice sites in a subregion is odd, since a well-defined Hilbert space of a Majorana fermion system must have an even number of Majorana fermions. We address the problem in this section and demonstrate how this even-odd problem, i.e., the \textit{fermion parity}, is intertwined with the entanglement structure. Understanding subsystem fermion parity provides crucial insights into formulating and resolving this ambiguity directly within the physical spin representation.

To resolve the issue of odd $N^A$, we add a redundant Majorana fermion, denoted by $c_{(N^A+1)}$, to the original system, making the number of Majorana fermions even, so that our previous approach works again~\cite{Lee-Wilczek_2013}. We design the redundant Majorana fermion in such a way that it has no interaction with other Majorana fermions. Mathematically, we set $\sanglelr{ c_{(N^A+1)}c_i}=0\;\forall i\in A$. In this way, the dimension of the Hilbert space of the matter sector, including the redundant one, is $2^{(N^A+1)/2}$, and the RDM becomes
\begin{equation}\label{equ.rho-A-F-odd}
    \rho_A^{F} = \frac{1}{2^{\lr{N^A+1}/2}}\sum_{[\mathfrak{D}]\in A}\mathfrak{C}_{\mathfrak{D}}\mathcal{M}_\mathfrak{D}
\end{equation}
The redundant Majorana fermion does not enter the DM. Thus, the form of $\rho^F_A$ remains intact. Only the normalization coefficient is changed. With the redundant Majorana fermion, the original correlation matrix $\Gamma^{(c)}$ is correspondingly enlarged by one dimension to $\Tilde{\Gamma}^{(c)}$. According to the definition of $c_{(N^A+1)}$, we obtain
\begin{equation}\label{equ.tilde-Gamma-corner}
    \Tilde{\Gamma}^{(c)} = 
\lr{
\begin{array}{c|c}
    \Gamma^{(c)} & 0 \\
    \hline
    0 & 1/2
\end{array}
}.
\end{equation}
In the enlarged correlation matrix $\Tilde{\Gamma}^{(c)}$, the upper-left block is the original correlation matrix. The upper-right column and the lower-left row are zeros because the redundant Majorana fermion does not correlate with others. The element $1/2$ at the lower-right entry is given by the redundant Majorana fermion $\frac{1}{2}=\frac{1}{2}\sanglelr{c_{(N^A+1)}c_{(N^A+1)}}$. 

Due to the change in the dimension of the Hilbert space of the matter sector, the relation \eqref{equ.U-V-trace-relation} needs to be modified to the following
\begin{equation}\label{equ.U-V-trace-relation-odd}
    \left\{
\begin{aligned}
&\Tr_{A} \mathcal{S}^{(A)}_{\mathfrak{D}} = 2^{N^A-N^A_p}\delta_{\mathfrak{D},\varnothing} \\
& \Tr_{A,F} \mathcal{M}_{\mathfrak{D}} = 2^{(N^A+1)/2}\delta_{\mathfrak{D},\varnothing}
\end{aligned}
\right..
\end{equation}
Performing the same calculation as we did in the cases with even $N^A$, we obtain a result similar to the EE 
\begin{align}\label{equ.EE-odd-derive-1}
    S_A &= -\Tr_{A,F} \rho^F_A\ln\rho^F_A + \lr{\frac{N^A-1}{2}-N^A_p}\ln 2 \nonumber\\
    &= -\tr\lr{\Tilde{\Gamma}^{(c)}\ln \Tilde{\Gamma}^{(c)}} + \lr{\frac{1}{2}\abs{\de A} - \frac{3}{2}}\ln 2,
\end{align}
where the correlation matrix in Eq.~\eqref{equ.EE-corr-mat} is replaced by the enlarged one. According to Eq.~\eqref{equ.tilde-Gamma-corner}, the first term of Eq.~\eqref{equ.EE-odd-derive-1} can be reduced and expressed in terms of the original correlation matrix 
\begin{equation}\label{equ.EE-reduce}
    \tr\lr{\Tilde{\Gamma}^{(c)} \ln \Tilde{\Gamma}^{(c)}} = \tr\lr{\Gamma^{(c)} \ln \Gamma^{(c)}} - \frac{1}{2}\ln 2 .
\end{equation}
Substituting Eq.~\eqref{equ.EE-reduce} into Eq.~\eqref{equ.EE-odd-derive-1}, the last term $-\frac{1}{2}\ln 2$ in Eq.~\eqref{equ.EE-reduce} will cancel the term $-\frac{1}{2}\ln 2$ in Eq.~\eqref{equ.EE-odd-derive-1} arising from the enlarged dimension of the matter sector, and we ultimately obtain a result that is the same as Eq.~\eqref{equ.Yao-Qi-standard-form} for even $N^A$. 

The above derivation not only confirms that the result in Ref.~\cite{Yao-Qi_2010} remains valid for subsystems with an odd number of sites $N^A$, but also provides a new physical insight into the interplay between the matter and gauge sectors. Operationally, we first borrow a dimension $2^{1/2}$ of a single Majorana fermion from the gauge sector to make the Hilbert space of the matter sector well-defined. The resulting shift of $-\frac{1}{2}\ln 2$ reflects this borrowed dimension. Subsequently, the redundant fermion yields the term $+\frac{1}{2}\ln 2$, effectively restoring the dimension back to the gauge sector and leaving no net change. 

Having established the entropy formula for odd $N^A$, we now examine the ES in this situation and its intimate connection with the subsystem fermion parity. As shown in Eq.~\eqref{equ.gamma-eigenvalue}, the elements in the single-particle ES come in pairs. Since the lower-right corner of Eq.~\eqref{equ.tilde-Gamma-corner} provides an eigenvalue $1/2$, there must be another identical eigenvalue $\frac{1}{2}=1-\frac{1}{2}$ in the original single-particle ES. Therefore, the many-body ES can be written as 
\begin{equation}\label{equ.many-body-entanglement-spectrum-odd}
\lambda_{n_1,n_2,\cdots,n_{(N^A+1)/2}} = \frac{1}{2}\prod_{k=1}^{\frac{N^A-1}{2}}\gamma_{k,n_k}
\end{equation}
The formula shows that the many-body ES depends on only $(N^A-1)/2$ quantum numbers, although the complete labeling contains $(N^A+1)/2$ of them. In other words, there is a twofold degeneracy in the ES of the matter sector. Thus, we only have $2^{\frac{N^A-1}{2}}$ distinct values of $\lambda_{\lrg{n}}$. The degree of degeneracy for each distinct eigenvalue is therefore $2^{\frac{N^A+1}{2}-N^A_p}$ because the total number of eigenvalues is equal to the number of distinct eigenvalues multiplied by the number of degeneracies
\[
2^{N^A-N^A_p} = 2^{\frac{N^A-1}{2}}\times 2^{\frac{N^A+1}{2}-N^A_p}.
\]
Due to the conservation of fermion parity in the subregion $A$, the fermion parity operator, defined by $\hat{\bog{\pi}}^{(A)} = i^{N^{A}}\prod_{j\in A}c_j$, commutes with the RDM. From the fermion-parity point of view, the twofold degeneracy in the matter sector can be understood as a block-diagonalization that depends on the subsystem fermion parity. As we have seen, the value of $\lambda_{\lrg{n}}$ is independent of the last fermion number $n_{(N^A+1)/2}$. At the same time, $\hat{\bog{\pi}}^{(A)}$ commutes with the entanglement Hamiltonian, and hence also commutes with $\rho_A$. Therefore, ES divides into even and odd sectors, and these two identical sectors correspond to $\lrg{n}$ with $n_{(N^A+1)/2}=0,1$, respectively, giving a twofold degeneracy in ES of the matter sector. The microscopic origin of this degeneracy in terms of the physical spin RDM is discussed in Sec.~\ref{sec:origin_degeneracy_odd}.

